\chapter{Instalação}

\section{Requisitos}

\hayashi{} é distribuído como binário estático para Linux, macOS e Windows. Não há
dependências de runtime: o executável \hay{} é autossuficiente.

\begin{itemize}
  \item Linux x86-64 ou ARM64
  \item macOS 12+ (Intel ou Apple Silicon)
  \item Windows 10+ (x86-64)
\end{itemize}

\section{Via Cargo (recomendado)}

Se Rust está instalado:

\begin{lstlisting}[style=inline, language=sh]
cargo install hayashi-lang
\end{lstlisting}

O binário \texttt{hay} será instalado em \texttt{\$HOME/.cargo/bin}. Confirme que
esse diretório está no \texttt{PATH}.

\begin{lstlisting}[style=inline, language=sh]
hay --version
\end{lstlisting}

% \section{Via AUR (Arch Linux)}
%
% \begin{lstlisting}[style=inline, language=sh]
% yay -S hayashi-lang
% \end{lstlisting}

\section{Binários pré-compilados}

Disponíveis na página de releases do repositório:
\url{https://github.com/sheep-farm/hayashi/releases}

Baixe o arquivo correspondente à sua plataforma, descompacte e coloque o binário
\texttt{hay} em algum diretório do \texttt{PATH}.

\section{Compilação a partir do código-fonte}

\begin{lstlisting}[style=inline, language=sh]
git clone https://github.com/sheep-farm/hayashi
cd hayashi
cargo build --release
# binário em target/release/hay
\end{lstlisting}

\section{Verificando a instalação}

\begin{lstlisting}[style=inline, language=sh]
hay --version
hay --help
\end{lstlisting}

Para executar um script:

\begin{lstlisting}[style=inline, language=sh]
hay meu_script.hay
\end{lstlisting}

Para entrar no REPL interativo:

\begin{lstlisting}[style=inline, language=sh]
hay
\end{lstlisting}

\begin{nota}
  No Arch Linux com Omarchy/Hyprland, \texttt{\$HOME/.cargo/bin} já costuma estar no
  \texttt{PATH} se o shell foi configurado com \texttt{rustup}. Caso contrário, adicione
  ao \texttt{.bashrc} ou \texttt{.zshrc}:
  \texttt{export PATH="\$HOME/.cargo/bin:\$PATH"}
\end{nota}
